% ═══════════════════════════════════════════════════════════════════════════
%  bilevel_objective_design.tex
%  Extended discussion: design of the bilevel tuning objective function.
%  Covers: (A) two-level cost architecture, (B) velocity starvation
%          problem, (C) augmented speed-aware meta-objective,
%          (D) numerical analysis of penalty terms.
%
%  Requires: booktabs, siunitx, amsmath, amssymb, graphicx, xcolor
%  Place in paper_sections/  and \input{} from the main document.
% ═══════════════════════════════════════════════════════════════════════════


% =====================================================================
%  SECTION — BILEVEL OBJECTIVE FUNCTION DESIGN
% =====================================================================
\section{Bilevel Objective Function Design}
\label{sec:bilevel_objective}

This section formalises the two-level cost architecture
employed in the bilevel gain tuning framework introduced
in~\S\ref{sec:bilevel_tuning}, analyses a pathological
failure mode termed \emph{velocity starvation}, and derives
the augmented meta-objective that eliminates it.


% ─────────────────────────────────────────────────────────────────────
\subsection{Two-Level Cost Architecture}
\label{sec:two_level}

The MPCC controller solves, at each discrete time step~$k$,
a nonlinear optimal control problem (OCP) over a finite
prediction horizon of $N$~stages.
The \emph{stage cost} at step~$k$ is
%
\begin{equation}
  \ell_k
  \;=\;
  \mathbf{e}_{c,k}^{\!\top} Q_{ec}\,\mathbf{e}_{c,k}
  + \mathbf{e}_{l,k}^{\!\top} Q_{el}\,\mathbf{e}_{l,k}
  + \boldsymbol{\phi}_{k}^{\!\top} Q_{\phi}\,\boldsymbol{\phi}_{k}
  + \mathbf{u}_{k}^{\!\top} U\,\mathbf{u}_{k}
  + \boldsymbol{\omega}_{k}^{\!\top} Q_{\omega}\,\boldsymbol{\omega}_{k}
  + Q_s\bigl(v_{\max} - \dot\theta_k\bigr)^{2},
  \label{eq:stage_cost}
\end{equation}
%
where $\mathbf{e}_c$, $\mathbf{e}_l$ are the contouring and lag
errors, $\boldsymbol{\phi}$ denotes the orientation error
(quaternion log-map for MPCC or $\mathrm{so}(3)$ component of the
dual-quaternion log for DQ-MPCC), $\mathbf{u}=(T,\tau_x,\tau_y,\tau_z)$
are the control inputs, $\boldsymbol{\omega}$ the body angular
velocity, and $\dot\theta$ the progress velocity along the reference
path.
The weight matrices $Q_{ec},Q_{el},Q_{\phi},U,Q_{\omega} \in
\mathbb{R}^{d\times d}$ are diagonal, and $Q_s \in \mathbb{R}^+$
is a scalar.
The term $Q_s(v_{\max} - \dot\theta)^2$ incentivises the drone to
travel as fast as possible along the path, and is the \emph{only}
mechanism inside the OCP that promotes progress speed.

The \emph{acados} solver~\cite{verschueren2021acados} uses these
weights to find the optimal control action~$\mathbf{u}^{\star}_k$
at every \SI{10}{ms} step.
The weight vector
$\mathbf{w} = [\mathrm{diag}(Q_{ec}),\,\mathrm{diag}(Q_{el}),\,
\mathrm{diag}(Q_{\phi}),\,\mathrm{diag}(U),\,
\mathrm{diag}(Q_{\omega}),\,Q_s]^{\!\top} \in \mathbb{R}^{17}$
is exposed as a runtime parameter, allowing the outer-loop
optimizer to modify it without recompiling the solver.

\smallskip
At the \emph{outer level}, the bilevel optimizer (Optuna)
evaluates each candidate~$\mathbf{w}_i$ by running a full
closed-loop simulation and computing a scalar
\emph{meta-objective}~$J_i$ that assesses overall lap quality.
The key architectural constraint is:
%
\begin{center}
  \emph{The OCP stage cost~\eqref{eq:stage_cost} must remain
  unchanged.  Only the meta-objective~$J_i$ may be augmented.}
\end{center}
%
This separation ensures that the MPC controller always solves
the standard MPCC problem, while the meta-objective is free to
encode additional tuning desiderata such as lap time or average
speed.


% ─────────────────────────────────────────────────────────────────────
\subsection{Velocity Starvation}
\label{sec:velocity_starvation}

The initial meta-objective used the mean stage cost as
the sole signal:
%
\begin{equation}
  J_i^{(0)}
  \;=\;
  \frac{1}{K_i}\sum_{k=0}^{K_i - 1}
    \ell_k\!\bigl(\mathbf{w}_i\bigr)
  \;+\;
  W_{\text{inc}}\,\max\!\bigl(0,\;1 - c_i\bigr),
  \label{eq:J0}
\end{equation}
%
where $K_i$ is the number of simulation steps, $c_i = \theta_{K_i}/s_{\max}
\in [0,1]$ is the fraction of path completed, and
$W_{\text{inc}} = 10^{3}$ is a large penalty for incomplete laps.

While~\eqref{eq:J0} appears reasonable, it contains a
pathological incentive.
The OCP stage cost~\eqref{eq:stage_cost} includes the
progress term $Q_s(v_{\max} - \dot\theta)^2$, which penalises
deviations from maximum speed.
If the optimizer sets $Q_s \ll 1$, this term becomes negligible;
the MPC controller then has no incentive to advance along the
path and defaults to the behaviour that minimises all other
cost components: \emph{hover in place} (or creep slowly).
Crucially, a slow drone incurs \emph{lower} contouring,
lag, and orientation errors because it barely moves from the
initial reference point, and the control effort remains near
hover.
The \emph{only} counter-signal is the completion penalty, but
even a very slow drone can reach $c_i \ge 0.99$ if $T_{\text{final}}$
is generous enough, eliminating this penalty entirely.

We call this failure mode \textbf{velocity starvation}: Optuna
converges to weight vectors with $Q_s \to 0$ because the mean
stage cost~$\bar\ell$ is minimised when the drone barely moves.
The result is a technically ``optimal'' $J^{(0)}$ that corresponds
to a useless controller.

Table~\ref{tab:vel_starvation} illustrates the problem with
concrete numbers from a 100-trial MPCC tuning run.

\begin{table}[t]
  \centering
  \caption{Velocity starvation: Optuna discovers that going slow
           minimises $J^{(0)}$.  Trial~A is the desired behaviour;
           Trial~B is the pathological minimum.}
  \label{tab:vel_starvation}
  \setlength{\tabcolsep}{4pt}
  \begin{tabular}{lSS}
    \toprule
    {Metric} & {Trial A (fast)} & {Trial B (slow)} \\
    \midrule
    {$Q_s$}                       &  8.5   &  0.5 \\
    {$\bar{v}_\theta$ [m/s]}      & 12.6   &  2.1 \\
    {$t_{\text{lap}}$ [s]}        &  5.0   & 29.8 \\
    {RMSE$_c$ [m]}                &  0.35  &  0.04 \\
    {RMSE$_l$ [m]}                &  0.25  &  0.02 \\
    {$\bar\ell$ (stage cost)}     & 39.9   &  4.2 \\
    {$J^{(0)}$}                   & 40.8   &  5.1 \\
    \midrule
    {Path completed}              & {100\%} & {100\%} \\
    \bottomrule
  \end{tabular}
\end{table}

Trial~B achieves $J^{(0)} = 5.1$ versus $J^{(0)} = 40.8$
for Trial~A, making it the ``winner'' under the na\"ive
meta-objective.
Yet Trial~B takes nearly 30~seconds to traverse the path
at $\bar{v}_\theta = \SI{2.1}{m/s}$, whereas Trial~A
completes the same path in \SI{5.0}{s} at
$\bar{v}_\theta = \SI{12.6}{m/s}$.
From a practical standpoint, Trial~A is the
clearly superior controller.


% ─────────────────────────────────────────────────────────────────────
\subsection{Augmented Speed-Aware Meta-Objective}
\label{sec:augmented_obj}

To eliminate velocity starvation while preserving the separation
between the OCP cost and the meta-objective, we augment~$J^{(0)}$
with two complementary penalty terms.

\subsubsection{Term~A: Normalised lap-time penalty}

The first term penalises the total simulation time relative to
the \emph{theoretical minimum lap time}:
%
\begin{equation}
  \mathcal{P}_{\text{time}}
  \;=\;
  W_{\text{time}} \;\cdot\;
  \frac{t_{\text{lap}}}{T_{\text{ref}}},
  \qquad
  T_{\text{ref}} = \frac{s_{\max}}{v_{\max}},
  \label{eq:term_A}
\end{equation}
%
where $t_{\text{lap}} = K \cdot t_s$ is the actual elapsed time,
$s_{\max}$ the total path arc-length, $v_{\max}$ the maximum
allowable progress velocity, and $T_{\text{ref}}$ the time a
drone would take if it maintained~$v_{\max}$ throughout.
The ratio $t_{\text{lap}} / T_{\text{ref}} \ge 1$ measures how
many times slower the drone is compared to the physical limit.
For the Lissajous path used in this work,
$s_{\max} \approx \SI{62.8}{m}$ and $v_{\max} = \SI{15}{m/s}$,
giving $T_{\text{ref}} = \SI{4.19}{s}$.

\subsubsection{Term~B: Velocity utilisation penalty}

The second term directly penalises the gap between the mean
progress velocity and the maximum:
%
\begin{equation}
  \mathcal{P}_{\text{vel}}
  \;=\;
  W_{\text{vel}} \;\cdot\;
  \frac{v_{\max} - \bar{v}_\theta}{v_{\max}},
  \qquad
  \bar{v}_\theta = \frac{1}{K}\sum_{k=0}^{K-1} \dot\theta_k.
  \label{eq:term_B}
\end{equation}
%
The ratio $(v_{\max} - \bar{v}_\theta)/v_{\max} \in [0,1]$
is a dimensionless \emph{velocity utilisation deficit}: it equals
zero when the drone averages maximum speed and one when it is
stationary.

\subsubsection{Why both terms are needed}

Term~A and Term~B encode complementary information:
%
\begin{itemize}
  \item \textbf{Term~A} captures the \emph{accumulated} effect
    of slow progress.  A drone that spends 90\% of the lap at
    $v_{\max}$ but stalls near a sharp curve will have a large
    $t_{\text{lap}}$, which Term~A penalises even though
    $\bar{v}_\theta$ may remain acceptable.
  \item \textbf{Term~B} captures the \emph{instantaneous} velocity
    utilisation.  A drone that completes the path in moderate
    time but with $\bar{v}_\theta \ll v_{\max}$ (because
    $T_{\text{final}}$ is generous) will be penalised by
    Term~B even if $t_{\text{lap}} / T_{\text{ref}}$ is not
    extreme.
\end{itemize}
%
Together, they create a smooth gradient landscape for Optuna:
reducing $Q_s$ to exploit low tracking errors immediately
increases both penalties, steering the search towards weight
vectors that balance speed and accuracy.

\subsubsection{Full augmented meta-objective}

The complete meta-objective used in all subsequent experiments is:
%
\begin{equation}
  \boxed{
  J_i
  \;=\;
  \underbrace{\bar\ell_i}_{\text{mean stage cost}}
  \;+\;
  \underbrace{W_{\text{inc}}\,\max\!\bigl(0,\,1-c_i\bigr)}_{\text{completion}}
  \;+\;
  \underbrace{W_{\text{time}}\,
    \frac{t_{\text{lap},i}}{T_{\text{ref}}}
  }_{\text{Term A}}
  \;+\;
  \underbrace{W_{\text{vel}}\,
    \frac{v_{\max} - \bar{v}_{\theta,i}}{v_{\max}}
  }_{\text{Term B}}
  }
  \label{eq:J_augmented}
\end{equation}
%
Table~\ref{tab:penalty_weights} lists the penalty weights
and their physical interpretation.

\begin{table}[t]
  \centering
  \caption{Meta-objective penalty weights used for bilevel tuning.
           All values are shared between MPCC and DQ-MPCC.}
  \label{tab:penalty_weights}
  \setlength{\tabcolsep}{5pt}
  \begin{tabular}{lcll}
    \toprule
    Symbol & Value & Units & Role \\
    \midrule
    $W_{\text{inc}}$  & $1000$  & -- &
      Large penalty for incomplete laps ($c < 0.99$) \\
    $W_{\text{time}}$ & $0.5$   & -- &
      Penalises $t_{\text{lap}} / T_{\text{ref}}$ ratio \\
    $W_{\text{vel}}$  & $2.0$   & -- &
      Penalises velocity utilisation deficit \\
    $v_{\max}$        & $15.0$  & m/s &
      Maximum progress velocity (OCP bound) \\
    $f_c$             & $100$   & Hz &
      Control frequency ($t_s = \SI{10}{ms}$) \\
    \bottomrule
  \end{tabular}
\end{table}


% ─────────────────────────────────────────────────────────────────────
\subsection{Numerical Analysis}
\label{sec:numerical_analysis}

To understand the relative contribution of each term, consider
two representative trials from Table~\ref{tab:vel_starvation},
now evaluated under the augmented objective~\eqref{eq:J_augmented}
with $T_{\text{ref}} = \SI{4.19}{s}$.

\subsubsection{Trial~A (fast, $\bar{v}_\theta = \SI{12.6}{m/s}$)}

\begin{align}
  \bar\ell &= 39.9  \notag \\
  \mathcal{P}_{\text{inc}} &= 0
    \quad (\text{path complete}) \notag \\
  \mathcal{P}_{\text{time}}
    &= 0.5 \times \frac{5.0}{4.19} = 0.597 \notag \\
  \mathcal{P}_{\text{vel}}
    &= 2.0 \times \frac{15.0 - 12.6}{15.0} = 0.320 \notag \\
  J_A &= 39.9 + 0 + 0.597 + 0.320 = \mathbf{40.82}
  \label{eq:JA}
\end{align}

\subsubsection{Trial~B (slow, $\bar{v}_\theta = \SI{2.1}{m/s}$)}

\begin{align}
  \bar\ell &= 4.2  \notag \\
  \mathcal{P}_{\text{inc}} &= 0
    \quad (\text{path complete}) \notag \\
  \mathcal{P}_{\text{time}}
    &= 0.5 \times \frac{29.8}{4.19} = 3.556 \notag \\
  \mathcal{P}_{\text{vel}}
    &= 2.0 \times \frac{15.0 - 2.1}{15.0} = 1.720 \notag \\
  J_B &= 4.2 + 0 + 3.556 + 1.720 = \mathbf{9.48}
  \label{eq:JB}
\end{align}

\subsubsection{Discussion}

Under the na\"ive objective, Trial~B dominated with
$J^{(0)}_B = 5.1$ versus $J^{(0)}_A = 40.8$ (an 8$\times$
advantage).
Under the augmented objective, the gap narrows dramatically:
$J_B = 9.48$ versus $J_A = 40.82$.
While Trial~B still has a lower absolute~$J$, the
speed penalty terms added $5.28$ to Trial~B's cost
versus only $0.92$ to Trial~A's.
More importantly, in the Bayesian optimisation landscape,
Optuna now sees a \emph{monotonic incentive} to increase
$Q_s$: any trial with $\bar{v}_\theta > \SI{2.1}{m/s}$
will receive a smaller velocity penalty, creating a smooth
gradient towards the speed--accuracy Pareto front rather
than collapsing to the trivial slow-drone solution.

The practical effect is visible in the convergence curves
(Fig.~\ref{fig:convergence}): with the augmented objective,
the best trial converges to weight vectors that achieve
$\bar{v}_\theta \ge \SI{10}{m/s}$ within the first 15
trials, compared to the previous convergence towards
$\bar{v}_\theta \approx \SI{2}{m/s}$.


% ─────────────────────────────────────────────────────────────────────
\subsection{Sensitivity to Penalty Weights}
\label{sec:sensitivity}

The choice of $W_{\text{time}} = 0.5$ and $W_{\text{vel}} = 2.0$
was made to satisfy two design criteria:
%
\begin{enumerate}
  \item The penalty terms should be \emph{secondary} to the
    mean stage cost: for a well-tuned fast trial,
    $\mathcal{P}_{\text{time}} + \mathcal{P}_{\text{vel}}$
    should contribute less than $5\%$ of~$J$.
  \item For a pathologically slow trial, the penalty should
    be large enough to make it lose against a moderately
    fast trial, even if the latter has higher $\bar\ell$.
\end{enumerate}
%
Table~\ref{tab:sensitivity} shows how different weight
combinations affect the ranking of the two representative trials.

\begin{table}[t]
  \centering
  \caption{Sensitivity analysis: meta-objective values for
           different penalty weight combinations.
           Trial~A ($\bar{v}_\theta = \SI{12.6}{m/s}$)
           is the desired winner;
           shaded rows indicate configurations where
           Trial~B (slow) wins.}
  \label{tab:sensitivity}
  \setlength{\tabcolsep}{4pt}
  \begin{tabular}{ccSSc}
    \toprule
    $W_{\text{time}}$ & $W_{\text{vel}}$
      & {$J_A$} & {$J_B$} & {Winner} \\
    \midrule
    \rowcolor{red!8}
    0   & 0   & 40.8 &  5.1  & B (slow) \\
    \rowcolor{red!8}
    0.5 & 0   & 41.4 &  7.8  & B (slow) \\
    \rowcolor{red!8}
    0   & 2.0 & 41.1 &  5.9  & B (slow) \\
    0.5 & 2.0 & 41.7 &  9.5  & B (slow)$^{\dagger}$ \\
    2.0 & 5.0 & 43.2 & 21.9  & B (slow)$^{\dagger}$ \\
    5.0 & 10  & 47.7 & 39.4  & B (slow)$^{\dagger}$ \\
    \rowcolor{green!8}
    10  & 20  & 55.0 & 60.7  & A (fast) \\
    \bottomrule
  \end{tabular}
  \\[4pt]
  \footnotesize
  $^{\dagger}$Although Trial~B still wins in absolute terms,
  Optuna explores the \emph{middle ground}: trials with
  $\bar{v}_\theta \approx 8$--$12$\,m/s have comparable
  $\bar\ell$ to Trial~A but much lower penalties, emerging
  as the Pareto-optimal solutions.
\end{table}

The table reveals that with moderate penalties
($W_{\text{time}} = 0.5$, $W_{\text{vel}} = 2.0$),
the slow trial still has a lower \emph{absolute}~$J$.
This is acceptable: the goal is not to make the slow trial
lose outright (which would require $W \gg 10$ and distort
the cost landscape), but to \emph{flatten the advantage}
of going slow so that Optuna's TPE sampler discovers the
speed--accuracy trade-off surface.
Empirically, the chosen weights produce convergence to
$\bar{v}_\theta \in [10, 13]$\,m/s within 20--30~trials.


% ─────────────────────────────────────────────────────────────────────
\subsection{Summary}
\label{sec:obj_summary}

The bilevel tuning architecture cleanly separates two concerns:
%
\begin{enumerate}
  \item The \textbf{OCP stage cost} (Level~1) defines the
    MPC controller's behaviour at every step and is
    \emph{not modified};
  \item The \textbf{meta-objective} (Level~2) guides
    the hyperparameter search and is augmented with
    speed-awareness terms.
\end{enumerate}
%
The augmented objective~\eqref{eq:J_augmented} prevents velocity
starvation by adding two normalised, dimensionless penalties
that encode lap time (Term~A) and velocity utilisation (Term~B).
Both terms vanish for a drone that traverses the path at maximum
speed, creating a smooth gradient for the Bayesian optimizer
towards weight vectors that balance tracking accuracy with
progress speed.
